\section{Main Results}
\label{sec:results}

We now prove the rank-count theorem. First, a linear-algebra lemma computes the dimension of the projected zero-violation space. Then a preservation lemma gives the geometric implication of zero labeled violation. Finally, the main theorem identifies the whole local solution code as an affine shift of that space.

\begin{lemma}[Branch-shift rank]
\label{lem:rank}
For the feasible branch shift space $\mathcal{K}_F = M(\ker V)$,
\[
\dim \mathcal{K}_F
=
\operatorname{rank}\begin{bmatrix} M \\ V \end{bmatrix}
-
\operatorname{rank}(V).
\]
\end{lemma}

\begin{proof}
Let $m = \lvert \mathcal{G} \rvert = \dim A$. Since $\mathcal{K}_F = M(\ker V)$, we restrict $M$ to $\ker V$:
\[
M\big|_{\ker V} : \ker V \to \mathbb{F}_2^{B_P}.
\]
By rank-nullity applied to this restricted map,
\[
\dim \mathcal{K}_F
=
\dim \ker V
-
\dim(\ker V \cap \ker M).
\]
Now $\dim \ker V = m - \operatorname{rank}(V)$, and
\[
\ker V \cap \ker M
=
\ker\begin{bmatrix} M \\ V \end{bmatrix}.
\]
Therefore,
\[
\dim(\ker V \cap \ker M)
=
m - \operatorname{rank}\begin{bmatrix} M \\ V \end{bmatrix}.
\]
Substituting these dimensions into the rank-nullity equation gives:
\[
\begin{aligned}
\dim \mathcal{K}_F &= (m - \operatorname{rank}(V)) - \left(m - \operatorname{rank}\begin{bmatrix} M \\ V \end{bmatrix}\right) \\
&= \operatorname{rank}\begin{bmatrix} M \\ V \end{bmatrix} - \operatorname{rank}(V).
\end{aligned}
\]
This completes the proof.
\end{proof}

The next lemma is the geometric direction encoded by the labeled violation matrix: zero net labeled violation gives a composition of mirror-compatible blocks, hence preserving every active edge.

\begin{lemma}[Zero-violation preservation]
\label{lem:preservation}
Let $\alpha \in \ker V$ and let $h = M\alpha$. If $s \in \Xi_F$, then $s \oplus h \in \Xi_F$.
\end{lemma}

\begin{proof}
We group the selected generators in $\alpha$ by mirror clique. For a fixed clique $C$, let $S_C$ be the XOR of the supports of the selected generators with mirror $C$. The row $(e,C)$ of $V\alpha$ represents the parity of selected generators with mirror $C$ that violate the active edge $e$. We have $V\alpha = 0$. Therefore, this parity is zero for every active edge $e \in F$ and every mirror clique $C$.

Therefore, the block $(S_C, C)$ is mirror-compatible with each active edge: if an active edge crosses $S_C$, its fixed endpoint lies in $C$. Generator supports are disjoint from their own mirror clique. Consequently, an endpoint in $C$ is never moved by a generator with mirror $C$. If a crossing edge had its fixed endpoint outside $C$, the selected generators with mirror $C$ would contribute an odd violation parity to the row $(e,C)$. This contradicts $V\alpha = 0$.

Reflection through the affine span $\operatorname{aff}(C)$ preserves edges with both endpoints moved, edges with neither endpoint moved, and crossing edges whose fixed endpoint lies directly on the mirror. Although spatial reflections generally do not commute, here any active edge crossing the block support $S_C$ has its fixed endpoint lying directly on the mirror $\operatorname{aff}(C)$. Consequently, each block reflection acts symmetrically on active edges, preserving the coordinates and affine hyperplanes representing subsequent mirror cliques. The composition of these reflections is therefore well-defined, and applying these block reflections one mirror clique at a time preserves every edge in $F$. Their net branch mask is $\bigoplus_C S_C = M\alpha = h$. Thus, $s \oplus h$ realizes every active edge in $F$.
\end{proof}

The main theorem states that, generically, the converse also holds: every feasible solution difference must come from such a zero-violation presentation.

\begin{theorem}[Feasible branch shift theorem]
\label{thm:main}
Assume $K \ge 2$, $E_D[L_P] \subseteq F$, $\Xi_F \ne \varnothing$, and generic mirror separation. Then, for any reference solution $s^\ast \in \Xi_F$,
\[
\Xi_F
=
s^\ast + \mathcal{K}_F
=
s^\ast + M(\ker V).
\]
\end{theorem}

\begin{proof}
Fix a reference solution $s^\ast \in \Xi_F$. We prove that the feasible branch differences from $s^\ast$ are exactly $\mathcal{K}_F$.

First, let $h \in \mathcal{K}_F$. By definition, $h = M\alpha$ for some $\alpha \in \ker V$. By Lemma~\ref{lem:preservation} (Zero-violation preservation),
\[
s^\ast \oplus h \in \Xi_F.
\]
Hence,
\[
s^\ast + \mathcal{K}_F \subseteq \Xi_F.
\]

Conversely, let $s \in \Xi_F$ and set $h = s \oplus s^\ast$. The cone generators form a triangular system spanning the full local branch space $\mathbb{F}_2^{B_P}$. Therefore, the branch shift $h$ has at least one generator presentation. If $h \notin \mathcal{K}_F$, then by Definition~\ref{def:mirror-sep} (Generic mirror separation), this branch shift changes at least one active-edge length outside a proper algebraic exceptional set. 

The coordinate parameters of the realization are generic; that is, they lie outside the union of the algebraic hypersurfaces where specific Cayley-Menger determinants or coordinate discriminants vanish. Consequently, any branch shift outside $\mathcal{K}_F$ must physically change the length of at least one active edge. However, both $s^\ast$ and $s$ realize every edge in $F$, meaning they preserve all active-edge lengths. This contradiction implies that $h$ must belong to $\mathcal{K}_F$, which yields
\[
\Xi_F \subseteq s^\ast + \mathcal{K}_F.
\]

Combining both inclusions yields:
\[
\Xi_F = s^\ast + \mathcal{K}_F = s^\ast + M(\ker V),
\]
which completes the proof.
\end{proof}

\begin{lemma}[Rank count]
\label{lem:count}
Under the assumptions of Theorem~\ref{thm:main},
\[
|\Xi_F|
=
2^{
\operatorname{rank}\begin{bmatrix} M \\ V \end{bmatrix}
-
\operatorname{rank}(V)
}.
\]
\end{lemma}

\begin{proof}
By Theorem~\ref{thm:main}, $\Xi_F$ is an affine coset of $\mathcal{K}_F$, so $|\Xi_F| = 2^{\dim \mathcal{K}_F}$. Applying Lemma~\ref{lem:rank} (Branch-shift rank) gives the displayed exponent.
\end{proof}
